% GENERATED FILE. Edit canonical lessons, then run npm run generate:books.
% canonical-source-hash: fnv1a64:a55ca7aad7c5eae0
% canonical-lessons: GU-C446-dhan, GU-C446-bharvad, GU-C446-goval, GU-C446-nes, GU-C446-caravvu, GU-R446-review-440-446, GU-R446-review-final-446

\chapter{Herds and Herders}
\label{ch:gu-herds-and-herders}
\hlchaptermodality{446}

\begin{chapteropening}
\textbf{By the end of this chapter:} \emph{I can talk about herders and their animals.}

Complete the last lesson of chapter 446: I can talk about herders and their animals.
\end{chapteropening}

\section[\texorpdfstring{\gu{ધણ} (dhaṇ)}{ધણ (dhaṇ)}]{\gu{ધણ} (dhaṇ) — a herd (of cattle)}

\label{lesson:GU-C446-dhan}



\begin{quote}
Before the new one: say the Gujarati for to call (of a bird), then the Gujarati for to chirp.
\end{quote}

\subsection*{You'll want to know: \gu{ધણ}}
\textbf{\gu{ધણ}} — \emph{dhaṇ} — \textquotedblleft{}a herd (of cattle)\textquotedblright{}.

\begin{quote}
\gu{સાંજે} \gu{ગાયોનું} \gu{ધણ} \gu{ગામમાં} \gu{પાછું} \gu{આવે} \gu{છે}.
\end{quote}

\begin{quote}
\emph{sā̃je gāyonũ dhaṇ gāmmā̃ pāchũ āve che.}
\end{quote}

\begin{quote}
In the evening the herd of cows comes back to the village.
\end{quote}

A neuter noun.

\begin{cousinweb}
\textbf{\gu{ધણ}} (\emph{dhaṇ}) is a tadbhava of Sanskrit \emph{dhana}, \textquotedblleft{}wealth\textquotedblright{}; CDIAL lists the Gujarati word, \textquotedblleft{}a herd of cattle\textquotedblright{} (CDIAL 6717): cattle were a family's wealth.
\end{cousinweb}

\begin{grammarlens}[title={What you've built}]
One more word for getting around.
\end{grammarlens}

\subsection*{Your turn}
\begin{itemize}
\raggedright
  \item \emph{Say it:} \emph{dhaṇ}
  \item \emph{Say it:} \emph{dhaṇ}, once more
  \item \emph{Say it:} \emph{sā̃je gāyonũ dhaṇ gāmmā̃ pāchũ āve che.} — In the evening the herd of cows comes back to the village.
  \item \emph{Recall:} say the Gujarati for to call (of a bird), then the Gujarati for to chirp, then say \emph{dhaṇ} again
\end{itemize}

\begin{tcolorbox}[breakable,colback=teal!4,colframe=teal!35!black,title={Before you move on}]
What does \gu{ધણ} mean? (\textquotedblleft{}A herd (of cattle)\textquotedblright{}.) Say it once more.
\end{tcolorbox}
\section[\texorpdfstring{\gu{ભરવાડ} (bharvāḍ)}{ભરવાડ (bharvāḍ)}]{\gu{ભરવાડ} (bharvāḍ) — a shepherd (of the Bharvad herders)}

\label{lesson:GU-C446-bharvad}



\begin{quote}
Before the new one: say the Gujarati for to chirp, then the Gujarati for a herd.
\end{quote}

\subsection*{You'll want to know: \gu{ભરવાડ}}
\textbf{\gu{ભરવાડ}} — \emph{bharvāḍ} — \textquotedblleft{}a shepherd (of the Bharvad herders)\textquotedblright{}.

\begin{quote}
\gu{ભરવાડ} \gu{ઘેટાં} \gu{લઈને} \gu{ડુંગર} \gu{પર} \gu{ગયો}.
\end{quote}

\begin{quote}
\emph{bharvāḍ gheṭā̃ laīne ḍungar par gayo.}
\end{quote}

\begin{quote}
The shepherd took the sheep up the hill.
\end{quote}

A masculine noun.

\begin{cousinweb}
\textbf{\gu{ભરવાડ}} (\emph{bharvāḍ}) goes back to \emph{*bhēḍrapāla}, \textquotedblleft{}a shepherd\textquotedblright{}, a form CDIAL reconstructs from \emph{bhēḍra}, a sheep, and \emph{pāla}, a keeper; CDIAL lists the Gujarati word, \textquotedblleft{}a shepherd or goatherd\textquotedblright{} (CDIAL 9608).
\end{cousinweb}

\begin{grammarlens}[title={What you've built}]
One more word for getting around.
\end{grammarlens}

\subsection*{Your turn}
\begin{itemize}
\raggedright
  \item \emph{Say it:} \emph{bharvāḍ}
  \item \emph{Say it:} \emph{bharvāḍ}, once more
  \item \emph{Say it:} \emph{bharvāḍ gheṭā̃ laīne ḍungar par gayo.} — The shepherd took the sheep up the hill.
  \item \emph{Recall:} say the Gujarati for to chirp, then the Gujarati for a herd, then say \emph{bharvāḍ} again
\end{itemize}

\begin{tcolorbox}[breakable,colback=teal!4,colframe=teal!35!black,title={Before you move on}]
What does \gu{ભરવાડ} mean? (\textquotedblleft{}A shepherd (of the Bharvad herders)\textquotedblright{}.) Say it once more.
\end{tcolorbox}
\section[\texorpdfstring{\gu{ગોવાળ} (govāḷ)}{ગોવાળ (govāḷ)}]{\gu{ગોવાળ} (govāḷ) — a cowherd}

\label{lesson:GU-C446-goval}



\begin{quote}
Before the new one: say the Gujarati for a herd, then the Gujarati for a shepherd.
\end{quote}

\subsection*{You'll want to know: \gu{ગોવાળ}}
\textbf{\gu{ગોવાળ}} — \emph{govāḷ} — \textquotedblleft{}a cowherd\textquotedblright{}.

\begin{quote}
\gu{ગોવાળ} \gu{ગાયો} \gu{સાથે} \gu{નદી} \gu{પાસે} \gu{બેઠો} \gu{હતો}.
\end{quote}

\begin{quote}
\emph{govāḷ gāyo sāthe nadī pāse beṭho hato.}
\end{quote}

\begin{quote}
The cowherd was sitting by the river with the cows.
\end{quote}

A masculine noun.

\begin{cousinweb}
\textbf{\gu{ગોવાળ}} (\emph{govāḷ}) is a tadbhava of Sanskrit \emph{gōpāla}, \textquotedblleft{}a protector of cows\textquotedblright{}, made of \emph{gō}, \textquotedblleft{}a cow\textquotedblright{}, and \emph{pāla}, \textquotedblleft{}a keeper\textquotedblright{} (CDIAL 4293).
\end{cousinweb}

\begin{grammarlens}[title={What you've built}]
One more word for getting around.
\end{grammarlens}

\subsection*{Your turn}
\begin{itemize}
\raggedright
  \item \emph{Say it:} \emph{govāḷ}
  \item \emph{Say it:} \emph{govāḷ}, once more
  \item \emph{Say it:} \emph{govāḷ gāyo sāthe nadī pāse beṭho hato.} — The cowherd was sitting by the river with the cows.
  \item \emph{Recall:} say the Gujarati for a herd, then the Gujarati for a shepherd, then say \emph{govāḷ} again
\end{itemize}

\begin{tcolorbox}[breakable,colback=teal!4,colframe=teal!35!black,title={Before you move on}]
What does \gu{ગોવાળ} mean? (\textquotedblleft{}A cowherd\textquotedblright{}.) Say it once more.
\end{tcolorbox}
\section[\texorpdfstring{\gu{નેસ} (nes)}{નેસ (nes)}]{\gu{નેસ} (nes) — a nes (a herders' hamlet in the forest)}

\label{lesson:GU-C446-nes}



\begin{quote}
Before the new one: say the Gujarati for a shepherd, then the Gujarati for a cowherd.
\end{quote}

\subsection*{You'll want to know: \gu{નેસ}}
\textbf{\gu{નેસ}} — \emph{nes} — \textquotedblleft{}a nes (a herders' hamlet in the forest)\textquotedblright{}.

\begin{quote}
\gu{જંગલમાં} \gu{ભરવાડોનો} \gu{નાનો} \gu{નેસ} \gu{છે}.
\end{quote}

\begin{quote}
\emph{jangalmā̃ bharvāḍono nāno nes che.}
\end{quote}

\begin{quote}
There is a small herders' hamlet in the forest.
\end{quote}

A masculine noun.

\begin{cousinweb}
\textbf{\gu{નેસ}} (\emph{nes}) is a tadbhava of Sanskrit \emph{nivēśa}, \textquotedblleft{}a dwelling place, a camp\textquotedblright{}; CDIAL lists the Gujarati word, \textquotedblleft{}a shepherd's hut\textquotedblright{} (CDIAL 7424).
\end{cousinweb}

\begin{grammarlens}[title={What you've built}]
One more word for getting around.
\end{grammarlens}

\subsection*{Your turn}
\begin{itemize}
\raggedright
  \item \emph{Say it:} \emph{nes}
  \item \emph{Say it:} \emph{nes}, once more
  \item \emph{Say it:} \emph{jangalmā̃ bharvāḍono nāno nes che.} — There is a small herders' hamlet in the forest.
  \item \emph{Recall:} say the Gujarati for a shepherd, then the Gujarati for a cowherd, then say \emph{nes} again
\end{itemize}

\begin{tcolorbox}[breakable,colback=teal!4,colframe=teal!35!black,title={Before you move on}]
What does \gu{નેસ} mean? (\textquotedblleft{}A nes (a herders' hamlet in the forest)\textquotedblright{}.) Say it once more.
\end{tcolorbox}
\section[\texorpdfstring{\gu{ચરાવવું} (carāvvũ)}{ચરાવવું (carāvvũ)}]{\gu{ચરાવવું} (carāvvũ) — to take to graze (animals)}

\label{lesson:GU-C446-caravvu}



\begin{quote}
Before the new one: say the Gujarati for a cowherd, then the Gujarati for a herders' hamlet.
\end{quote}

\subsection*{You'll want to know: \gu{ચરાવવું}}
\textbf{\gu{ચરાવવું}} — \emph{carāvvũ} — \textquotedblleft{}to take to graze (animals)\textquotedblright{}.

\begin{quote}
\gu{ગોવાળ} \gu{રોજ} \gu{ગાયોને} \gu{મેદાનમાં} \gu{ચરાવે} \gu{છે}.
\end{quote}

\begin{quote}
\emph{govāḷ roj gāyone medānmā̃ carāve che.}
\end{quote}

\begin{quote}
Every day the cowherd takes the cows to graze in the field.
\end{quote}

Past, with \gu{તેણે} (he or she …): \gu{તેણે} \gu{ગાયો} \gu{ચરાવી}, \emph{teṇe gāyo carāvī}.

\begin{cousinweb}
\textbf{\gu{ચરાવવું}} (\emph{carāvvũ}) is the causative partner of \gu{ચરવું}, \textquotedblleft{}to graze\textquotedblright{}, a tadbhava of Sanskrit \emph{carati}, \textquotedblleft{}moves, grazes\textquotedblright{}; CDIAL lists the Gujarati verb (CDIAL 4686).
\end{cousinweb}

\begin{grammarlens}[title={What you've built}]
One more word for getting around.
\end{grammarlens}

\subsection*{Your turn}
\begin{itemize}
\raggedright
  \item \emph{Say it:} \emph{carāvvũ}
  \item \emph{Say it:} \emph{carāvvũ}, once more
  \item \emph{Say it:} \emph{govāḷ roj gāyone medānmā̃ carāve che.} — Every day the cowherd takes the cows to graze in the field.
  \item \emph{Recall:} say the Gujarati for a cowherd, then the Gujarati for a herders' hamlet, then say \emph{carāvvũ} again
\end{itemize}

\begin{tcolorbox}[breakable,colback=teal!4,colframe=teal!35!black,title={Before you move on}]
What does \gu{ચરાવવું} mean? (\textquotedblleft{}To take to graze (animals)\textquotedblright{}.) Say it once more.
\end{tcolorbox}
\section[(dialogue)]{Looking back: from Streams and Canals to Herds and Herders}

\label{lesson:GU-R446-review-440-446}



\begin{quote}
Say the Gujarati for a mountain stream and a current.
\end{quote}

\subsection*{Your turn}
Say these aloud:

\begin{itemize}
\raggedright
  \item a mountain stream, a current, a confluence, a canal, to let flow
  \item a large lake, a dam, a stepwell, a riverbank, to sink in
  \item the ocean, a wave, a river boat, a fishing net, to rock gently
  \item a tree, a banyan, a peepal, a mango tree, to sprout
  \item a tree trunk, bark, a bud, a cluster, to be fragrant
  \item a nest, a firefly, a cricket (the insect), to call (of a bird), to chirp
  \item a herd, a shepherd, a cowherd, a herders' hamlet, to take to graze
\end{itemize}

\begin{tcolorbox}[breakable,colback=teal!4,colframe=teal!35!black,title={Before you move on}]
A mountain stream? (\textbf{\gu{ઝરણું}}, \emph{jharṇũ}.) A current? (\textbf{\gu{પ્રવાહ}}, \emph{pravāh}.) A confluence? (\textbf{\gu{સંગમ}}, \emph{sangam}.) A canal? (\textbf{\gu{નહેર}}, \emph{naher}.) To let flow? (\textbf{\gu{વહાવવું}}, \emph{vahāvvũ}.) A large lake? (\textbf{\gu{સરોવર}}, \emph{sarovar}.) A dam? (\textbf{\gu{ડેમ}}, \emph{ḍem}.) A stepwell? (\textbf{\gu{વાવ}}, \emph{vāv}.) A riverbank? (\textbf{\gu{કાંઠો}}, \emph{kā̃ṭho}.) To sink in? (\textbf{\gu{ખૂંપવું}}, \emph{khū̃pvũ}.) The ocean? (\textbf{\gu{સમુદ્ર}}, \emph{samudra}.) A wave? (\textbf{\gu{લહેર}}, \emph{laher}.) A river boat? (\textbf{\gu{નાવ}}, \emph{nāv}.) A fishing net? (\textbf{\gu{જાળ}}, \emph{jāḷ}.) To rock gently? (\textbf{\gu{ડોલવું}}, \emph{ḍolvũ}.) A tree? (\textbf{\gu{ઝાડ}}, \emph{jhāḍ}.) A banyan? (\textbf{\gu{વડ}}, \emph{vaḍ}.) A peepal? (\textbf{\gu{પીપળો}}, \emph{pīpḷo}.) A mango tree? (\textbf{\gu{આંબો}}, \emph{ā̃bo}.) To sprout? (\textbf{\gu{ફૂટવું}}, \emph{phūṭvũ}.) A tree trunk? (\textbf{\gu{થડ}}, \emph{thaḍ}.) Bark? (\textbf{\gu{છાલ}}, \emph{chāl}.) A bud? (\textbf{\gu{કળી}}, \emph{kaḷī}.) A cluster? (\textbf{\gu{ઝૂમખું}}, \emph{jhūmkhũ}.) To be fragrant? (\textbf{\gu{મહેકવું}}, \emph{mahekvũ}.) A nest? (\textbf{\gu{માળો}}, \emph{māḷo}.) A firefly? (\textbf{\gu{આગિયો}}, \emph{āgiyo}.) A cricket (the insect)? (\textbf{\gu{તમરું}}, \emph{tamrũ}.) To call (of a bird)? (\textbf{\gu{ટહુકવું}}, \emph{ṭahukvũ}.) To chirp? (\textbf{\gu{ચહેકવું}}, \emph{cahekvũ}.) A herd? (\textbf{\gu{ધણ}}, \emph{dhaṇ}.) A shepherd? (\textbf{\gu{ભરવાડ}}, \emph{bharvāḍ}.) A cowherd? (\textbf{\gu{ગોવાળ}}, \emph{govāḷ}.) A herders' hamlet? (\textbf{\gu{નેસ}}, \emph{nes}.) To take to graze? (\textbf{\gu{ચરાવવું}}, \emph{carāvvũ}.)
\end{tcolorbox}
\section[(dialogue)]{Herds and Herders — once more}

\label{lesson:GU-R446-review-final-446}



\begin{quote}
Say the Gujarati for a herd and a shepherd.
\end{quote}

\subsection*{Your turn}
\begin{itemize}
\raggedright
  \item \emph{Say it:} a herd, a shepherd, a cowherd, a herders' hamlet, to take to graze
\end{itemize}

\begin{tcolorbox}[breakable,colback=teal!4,colframe=teal!35!black,title={Before you move on}]
A herd? (\textbf{\gu{ધણ}}, \emph{dhaṇ}.) A shepherd? (\textbf{\gu{ભરવાડ}}, \emph{bharvāḍ}.) A cowherd? (\textbf{\gu{ગોવાળ}}, \emph{govāḷ}.) A herders' hamlet? (\textbf{\gu{નેસ}}, \emph{nes}.) To take to graze? (\textbf{\gu{ચરાવવું}}, \emph{carāvvũ}.)
\end{tcolorbox}
